\begin{tikzpicture}[baseline=(current bounding box.center), x=0.85cm, y=0.6cm]
  % U = matriz de vecindades mínimas de G_{13}, lectura por filas como conjuntos
  % Resaltado de la fila que se está leyendo: {fila/diapositiva}
  \foreach \fil/\dia in {1/2, 2/3, 3/4, 4/5}
    \fill<\dia>[FancyBlue, opacity=0.25, rounded corners=2pt] (0.6,-\fil-0.38) rectangle (4.4,-\fil+0.38);
  \node at (-1.65,-2.5) {$U=$};
  % Etiquetas de filas y columnas
  \foreach \k in {1,...,4} {
    \node at (\k,0.1) {$v_{\k}$};
    \node at (-0.3,-\k) {$U_{v_{\k}}$};
  }
  % Paréntesis
  \draw[thick] (0.55,-0.55) to[bend right=25] (0.55,-4.45);
  \draw[thick] (4.45,-0.55) to[bend left=25] (4.45,-4.45);
  % Entradas: {fila/columna/valor}
  \foreach \fil/\col/\val in {1/1/1, 1/2/0, 1/3/0, 1/4/0,
                              2/1/0, 2/2/1, 2/3/0, 2/4/0,
                              3/1/0, 3/2/0, 3/3/1, 3/4/0,
                              4/1/1, 4/2/0, 4/3/1, 4/4/1}
    \node at (\col,-\fil) {$\val$};
  % Viñetas: cada vecindad mínima a la altura de su fila (overlay: no cambian el tamaño de la figura)
  \foreach \fil/\conj/\dia in {1/{\{v_1\}}/2, 2/{\{v_2\}}/3, 3/{\{v_3\}}/4, 4/{\{v_1,v_3,v_4\}}/5}
    \node<\dia->[anchor=west, overlay] at (5.6,-\fil)
      {{\usebeamercolor[fg]{itemize item}\usebeamertemplate{itemize item}}\ \ $U_{v_{\fil}}=\conj$};
\end{tikzpicture}
